\chapter{反应机理：弯箭头}\label{ch:g12:curly-arrows}

配平的方程式就像电影的第一张和最后一张照片：
它展示反应前的反应物和反应后的生成物，
而两者之间的过程一概不显示。然而，分子并不是简单地交换原子。
键是一次断裂或形成一个，电子对在不同位置之间移动，
经过一些寿命很短、从不出现在方程式中的物种。
化学家用弯箭头来表示这些移动。
掌握几条简单规则后，这些箭头便能解释反应为何在某处发生，
并预测从未见过的反应。

\begin{recall}
电负性、极性键和部分电荷
（\cref{ch:g11:polarity-and-cohesion}）。路易斯结构：成键电子对和
孤电子对（\cref{ch:g10:lewis-and-shape}）。官能团与有机物
类别（\cref{ch:g11:functional-groups}）。
\end{recall}

\section{富电子位点与缺电子位点}

\begin{definition}[给电子位点与受电子位点]\label{def:g12:curly-arrows:site}
分子或离子中的\emph{给电子位点}\index{给电子位点}是富电子的位置，
能够提供一对电子以形成新键，包括孤电子对、双键、
带负电荷的原子或
带部分负电荷（$\delta^-$）的原子。\emph{受电子位点}\index{受电子位点}是缺电子的位置，
能够接受一对电子，即带正电荷的原子或
带部分正电荷的原子
（$\delta^+$）。
\end{definition}

\begin{example}[两个分子中的位点]\label{ex:g12:curly-arrows:sites}
在溴乙烷 \ce{CH3-CH2-Br} 中，溴比碳的电负性大，与溴相连的碳带 $\delta^+$（受电子位点），
溴带 $\delta^-$，并有三对孤电子对。
在丙酮中，\ce{C=O} 基团的氧带 $\delta^-$，
有两对孤电子对（给电子位点），其碳带 $\delta^+$（受电子位点）。
氢氧根离子 \ce{OH-} 带负电荷、有三对孤电子对，
是较强的给电子体。
\end{example}

\begin{omfigure}
\setchemfig{atom sep=2.2em}
\begin{tabular}{cc}
\chemfig{H_3C-C(-[2]H)(-[6]H)-\charge{90:2pt=\:,0=\:,270:2pt=\:}{Br}} &
\chemfig{H_3C-C(-[6]CH_3)=[2]\charge{0=\:,180=\:}{O}} \\[2pt]
\small $\delta^+$ 位于与 \ce{Br} 相连的碳上，$\delta^-$ 位于 \ce{Br} 上 &
\small $\delta^+$ 位于 \ce{C=O} 的碳上，$\delta^-$ 位于 \ce{O} 上 \\
\small 溴乙烷 & \small 丙酮
\end{tabular}

{\small 缺电子的碳原子（$\delta^+$，受电子位点）与
富电子、电负性大的原子（$\delta^-$，带孤电子对，
即给电子位点）相邻。}
\end{omfigure}

\section{弯箭头}

\begin{definition}[弯箭头]\label{def:g12:curly-arrows:curly-arrow}
\emph{弯箭头}\index{弯箭头}表示反应某一步中一对电子的移动，
它始于移动的那一对电子，即孤电子对或化学键，
终于电子所去的位置：
既可以是在某原子上形成新键，
也可以是接受断裂之键的电子对的原子。
\end{definition}

\begin{method}[画弯箭头]\label{met:g12:curly-arrows:drawing}
\begin{enumerate}
  \item 找出给电子位点和受电子位点。
  \item 从给电子的电子对（孤电子对或化学键）画一个箭头指向
        受电子原子：该处形成新键。
  \item 若受电子原子此时的电子数过多（超过八电子，
        对氢原子则超过两个电子），再画一个箭头断开它的一个键，
        电子对归于电负性较大的
        原子。
  \item 检查该步之后的电荷：给出孤电子对的原子
        正电荷增加 1，接受电子对、形成孤电子对的原子
        负电荷增加 1；总电荷保持不变。
\end{enumerate}
\end{method}

\begin{proposition}[每一步中电荷守恒]\label{prop:g12:curly-arrows:charge}
物种的总电荷在
机理的每一步前后都相同。
\end{proposition}

\begin{proof}
弯箭头只移动原本就存在的电子；
既不产生也不消灭电子，原子核也不变化。
\end{proof}

\begin{example}[第一个箭头]\label{ex:g12:curly-arrows:first}
氢离子 \ce{H+} 与氢氧根离子 \ce{OH-} 相遇。氧的一对孤电子对
与没有电子的氢离子形成一个键：
\ce{H+ + OH- -> H2O}。画一个箭头，从氧的孤电子对指向
\ce{H+}；电荷 $+1$ 与 $-1$ 在水分子中变为 0。
\end{example}

\begin{omfigure}
\setchemfig{atom sep=2.1em}
\chemfig{@{hp}H^{+}}\qquad\raisebox{0.2ex}{$+$}\qquad
\chemfig{H-@{ox}\charge{90=\:,0=\:,270=\:}{O}^{-}}
\qquad\raisebox{0.2ex}{$\longrightarrow$}\qquad
\chemfig{H-\charge{90=\:,270=\:}{O}-H}
\chemmove{
  \draw[->, thick, omProp, shorten <=4pt, shorten >=3pt]
    ($(ox)+(0,0.25)$) .. controls +(110:9mm) and +(60:9mm) .. (hp);}

{\small 最简单的一步：氢氧根离子的一对孤电子对移动，以形成
与氢离子之间的新 \ce{O-H} 键。}
\end{omfigure}

\section{基本步骤与中间体}

\begin{definition}[机理、基元反应、中间体]\label{def:g12:curly-arrows:mechanism}
某一反应的\emph{反应机理}\index{反应机理}是
反应物转变为生成物所经历的一系列步骤。每一个
\emph{基元反应}\index{基元反应}是一次单独的事件，其中
若干对电子同时发生移动。在某一步中生成、又在后一步中被消耗的
物种称为\emph{反应中间体}\index{反应中间体}：
它不出现在总反应方程式中。
\end{definition}


\begin{example}[碳正离子中间体]\label{ex:g12:curly-arrows:carbocation}
溴化氢与乙烯发生加成时，反应分两步进行。
第一步中，乙烯的双键把一对电子给予
\ce{HBr} 中的氢原子，\ce{H-Br} 的电子对归于溴原子，溴以
\ce{Br-} 的形式离去；失去双键共用电子对的碳原子只剩
六个电子并带正电荷：碳正离子
\ce{CH3-CH2+}。第二步中，\ce{Br-} 的一对孤电子对与
该碳原子成键。碳正离子是一种中间体。
\end{example}

\section{三类反应}

\begin{definition}[取代、加成、消除]\label{def:g12:curly-arrows:categories}
在\emph{取代}\index{取代}反应中，分子中的一个原子或原子团
被另一原子或原子团所替代。在\emph{加成}\index{加成}反应中，
两个分子结合为一个分子，其中一个多重键变为单键。在
\emph{消除}\index{消除}反应中，分子中脱去一个小分子（通常是水），
单键变为多重键。
\end{definition}

\begin{omfigure}
\setchemfig{atom sep=2em}
\small
\textbf{取代}（一步完成）：\ce{CH3-CH2-Br + OH- -> CH3-CH2-OH + Br-}

\medskip
\chemfig{H-@{o1}\charge{90=\:,0=\:,270=\:}{O}^{-}}\qquad
\chemfig{H_3C-@{c1}C(-[2]H)(-[6]H)-[@{b1}]@{br1}\charge{90=\:,0=\:,270=\:}{Br}}
\qquad$\longrightarrow$\qquad
\chemfig{H_3C-C(-[2]H)(-[6]H)-O-H}\quad$+$\quad\ce{Br-}
\chemmove{
  \draw[->, thick, omProp, shorten <=4pt, shorten >=3pt]
    ($(o1)+(0.15,0.2)$) .. controls +(60:9mm) and +(120:9mm) .. (c1);
  \draw[->, thick, omDef, shorten <=2pt, shorten >=4pt]
    (b1) .. controls +(-90:7mm) and +(-90:7mm) .. (br1);}

\vspace{6mm}
\textbf{加成}（两步完成）：\ce{CH2=CH2 + HBr -> CH3-CH2-Br}

\vspace{9mm}
\chemfig{H_2C=[@{pi}]CH_2}\quad$+$\quad
\chemfig{@{h2}H-[@{b2}]@{br2}Br}
\qquad$\longrightarrow$\qquad
\chemfig{H_3C-@{cp}\charge{90:2pt=$\scriptstyle +$}{C}H_2}\quad$+$\quad
\chemfig{@{brm}\charge{90=\:,180=\:,270=\:,0=\:}{Br}^{-}}
\qquad$\longrightarrow$\qquad \ce{CH3-CH2-Br}
\chemmove{
  \draw[->, thick, omProp, shorten <=2pt, shorten >=3pt]
    (pi) .. controls +(90:10mm) and +(110:9mm) .. (h2);
  \draw[->, thick, omDef, shorten <=2pt, shorten >=4pt]
    (b2) .. controls +(-90:7mm) and +(-90:7mm) .. (br2);
  \draw[->, thick, omProp, shorten <=4pt, shorten >=3pt]
    ($(brm)+(0,-0.25)$) .. controls +(-120:8mm) and +(-60:8mm) .. (cp);}

\vspace{6mm}
\textbf{消除}（乙醇的酸催化脱水，三步完成）：
\ce{CH3-CH2-OH -> CH2=CH2 + H2O}

\medskip
\begin{tabular}{@{}l@{}}
1.\ 氧原子获得一个氢离子：\ce{CH3-CH2-OH + H+ -> CH3-CH2-OH2+}；\\
2.\ \ce{C-O} 电子对与水一起离去：\ce{CH3-CH2-OH2+ -> CH3-CH2+ + H2O}；\\
3.\ 相邻的 \ce{C-H} 电子对形成双键，同时
    \ce{H+} 离去：\\
\phantom{3.\ }\ce{CH3-CH2+ -> CH2=CH2 + H+}。
\end{tabular}

\medskip
{\small 三种机理。橙色箭头：给电子对形成新键。
蓝色箭头：键断裂，其电子对归于电负性更大的
原子。消除反应中，氢离子最后被重新释放。}
\end{omfigure}


\begin{remark}[改变骨架或官能团]\label{rem:g12:curly-arrows:skeleton}
有些反应只改变官能团，而保持碳骨架不变：
溴乙烷的取代反应将卤代烷转化为醇，
反应前后各含两个碳原子。另一些反应则构建或断裂
骨架：形成新的碳碳键会使碳链加长，这是
由小分子合成大分子的关键步骤。设计合成路线
需要按正确顺序选择这两类反应。
\end{remark}

\section{练习}

\begin{exercise}[$\star$]\label{exo:g12:curly-arrows:1}
在下列每种物种中，指出一个给电子位点，若存在，再指出一个
受电子位点：\ce{OH-}；\ce{H+}；\ce{CH2=CH2}；\ce{CH3-Cl}；\ce{H2O}。
\end{exercise}

\begin{exercise}[$\star$]\label{exo:g12:curly-arrows:2}
下列反应分别属于取代、加成还是消除反应？
(a) \ce{CH3-CH2-Cl + OH- -> CH3-CH2-OH + Cl-}；
(b) \ce{CH2=CH2 + H2O -> CH3-CH2-OH}；
(c) \ce{CH3-CH2-OH -> CH2=CH2 + H2O}；
(d) \ce{CH2=CH2 + Br2 -> CH2Br-CH2Br}。
\end{exercise}

\begin{exercise}[$\star$]\label{exo:g12:curly-arrows:3}
弯箭头从哪里开始，在哪里结束？
从氧原子的孤电子对指向碳原子的弯箭头表示什么？
\end{exercise}

\begin{exercise}[$\star$]\label{exo:g12:curly-arrows:4}
在反应 \ce{H+ + NH3 -> NH4+} 中，哪种物种提供电子对？
画出弯箭头。生成物带何种电荷？
\end{exercise}

\begin{exercise}[$\star$]\label{exo:g12:curly-arrows:5}
什么是反应中间体？请写出
\ce{HBr} 与乙烯加成反应的中间体。
\end{exercise}

\begin{exercise}[$\star\star$]\label{exo:g12:curly-arrows:6}
画出取代反应
\ce{CH3-Br + OH- -> CH3-OH + Br-}（一步完成）的弯箭头，并检查电荷。
\end{exercise}

\begin{exercise}[$\star\star$]\label{exo:g12:curly-arrows:7}
在 \ce{HBr} 与乙烯加成的第一步中，为什么
\ce{H-Br} 电子对归属于溴而不是氢？
该步每种产物分别带何种电荷？
\end{exercise}

\begin{exercise}[$\star\star$]\label{exo:g12:curly-arrows:8}
写出乙醇脱水的三步反应，并对每一步
指出哪个物种是给电子位点、哪个是受电子位点。
为什么氢离子不写入总反应方程式？
\end{exercise}

\begin{exercise}[$\star\star$]\label{exo:g12:curly-arrows:9}
步骤 \ce{CH3-CH2+ + H2O -> CH3-CH2-OH2+} 紧接在
乙烯于水中加成氢离子之后。画出弯箭头。
再失去 \ce{H+} 后生成什么？写出总反应方程式，并
说明其类别。
\end{exercise}

\begin{exercise}[$\star\star$]\label{exo:g12:curly-arrows:10}
阅读溴乙烷取代反应的机理：有几对
电子对发生移动？哪个键形成，哪个键断裂？是否有中间体？
\end{exercise}

\begin{exercise}[$\star\star$]\label{exo:g12:curly-arrows:11}
% ledger: en:C, en:Cl, en:Br
氯甲烷与氢氧根离子的反应比溴甲烷
慢。利用电负性（C 2.55，Cl 3.16，Br 2.96），
解释哪个碳原子的 $\delta^+$ 更大。这能否
解释二者速率的差异？（还应考虑离去离子夺取电子对的难易程度。）
\end{exercise}

\begin{exercise}[$\star\star\star$]\label{exo:g12:curly-arrows:12}
氯化氢加成到丙烯 \ce{CH2=CH-CH3} 上。第一步中
氢离子可与双键两端的任一碳原子成键，
生成两种可能的碳正离子。画出这两种碳正离子及两种可能的产物，并
为它们命名。（哪一种生成最多，在本科一年级分册中学习。）
\end{exercise}

\begin{exercise}[$\star\star\star$]\label{exo:g12:curly-arrows:13}
一名学生画出溴乙烷取代反应的第一步，
其弯箭头从溴原子指向 \ce{OH-} 中的氧原子。请解释
其中的两处错误。
\end{exercise}

\begin{exercise}[$\star\star\star$]\label{exo:g12:curly-arrows:14}
由羧酸与醇生成酯的过程中，醇的氧原子
进攻 \ce{C=O} 基团中的碳原子。利用
部分电荷，解释为什么该碳原子是受电子位点、
该氧原子是给电子位点。总反应失去的是哪种小分子，以及
整个反应属于哪一类？
\end{exercise}

\begin{exercise}[$\star\star\star$]\label{exo:g12:curly-arrows:15}
乙烯与溴 \ce{Br2}（非极性分子）反应。当 \ce{Br2}
分子接近双键时，\ce{Br-Br} 的电子被推离靠近的溴原子。
解释为何这会形成一个
受电子位点，并提出机理的第一步。
\end{exercise}

\section{问题：制备溴乙烷}

\begin{problem}[{周末问题——从 10 克乙烯出发：能
得到多少溴乙烷？}]
\label{pb:g12:curly-arrows:1}
% ledger: aw:C, aw:H, aw:Br, en:H, en:Br, en:C
溴乙烷是许多合成的起始原料，可通过
溴化氢与乙烯加成制得：\ce{CH2=CH2 + HBr -> CH3-CH2-Br}。一位化学家
用 \qty{10.0}{g} 乙烯和过量的溴化氢进行反应，
其产率为 \qty{75}{\percent}。

\medskip
\noindent\textbf{第一部分——位点。}
\begin{enumerate}
  \item 画出乙烯和溴化氢的路易斯结构。
  \item 乙烯的给电子位点在哪里？为什么？
  \item 计算 \ce{H-Br} 键的电负性之差
        （H 2.20，Br 2.96）。哪个原子带 $\delta^+$ 电荷？
  \item \ce{HBr} 中哪个原子是受电子位点？
\end{enumerate}

\medskip
\noindent\textbf{第二部分——反应机理。}
\begin{enumerate}[resume]
  \item 用两个弯箭头画出第一步反应。
  \item 写出生成的两种物种及其电荷。
  \item 检验第一步中电荷是否守恒。
  \item 用弯箭头画出第二步反应。
  \item 中间体中带正电的碳原子周围有多少个电子？
        这种中间体为何如此活泼？
\end{enumerate}

\medskip
\noindent\textbf{第三部分——属于哪类反应？}
\begin{enumerate}[resume]
  \item 该反应机理中的中间体是什么？
  \item 它为什么不出现在总反应方程式中？
  \item 该反应属于哪一类？
  \item 该反应改变了碳骨架，还是只改变
        官能团？
\end{enumerate}

\medskip
\noindent\textbf{第四部分——能得到多少产物？}
\begin{enumerate}[resume]
  \item 计算乙烯和溴乙烷的摩尔质量。
  \item 计算乙烯的物质的量。
  \item 填写进度表，哪种反应物是限量反应物？
  \item 计算溴乙烷的最大质量。
  \item 计算产率为 \qty{75}{\percent} 时得到的质量。
  \item 给出最终答案：化学家能得到
        多少质量的溴乙烷？
\end{enumerate}
\end{problem}
